\documentclass[11pt]{article}
\usepackage[a4paper,margin=22mm]{geometry}
\usepackage{fontspec}
\setmainfont{DejaVu Serif}
\setsansfont{DejaVu Sans}
\usepackage{microtype}
\usepackage{xcolor}
\usepackage{hyperref}
\usepackage{enumitem}
\usepackage{parskip}
\definecolor{Accent}{HTML}{971D32}
\definecolor{Muted}{HTML}{666666}
\hypersetup{colorlinks=true,urlcolor=Accent,linkcolor=Accent}
\setlist{leftmargin=1.6em,itemsep=0.3em,topsep=0.35em}
\setcounter{tocdepth}{2}
\renewcommand{\contentsname}{Contents}
\newcommand{\sectionrule}{\par\vspace{-0.5em}\noindent\textcolor{Accent}{\rule{\linewidth}{0.6pt}}\par}
\begin{document}

\begin{center}
{\sffamily\bfseries\color{Accent} TOEFL WORD OF THE DAY\par}
\vspace{8mm}
{\fontsize{42}{48}\selectfont\bfseries corresponding\par}
\vspace{2mm}
{\Large\sffamily\color{Accent} /ˌkɔːrəˈspɑːndɪŋ/\par}
\vspace{2mm}
{\small\sffamily Local audio phrase: corresponding\par}
{\small\sffamily Audio file: audios/corresponding.mp3\par}
\vspace{2mm}
{\large\sffamily\color{Muted} adjective\par}
\vspace{5mm}
{\sffamily directly related or matching \textbullet\ equivalent in another system or group\par}
\vspace{6mm}
{\large Corresponding describes things that are directly connected, match each other, or occupy equivalent positions.\par}
\vspace{6mm}
\begin{minipage}{0.84\textwidth}
\itshape “An increase in production was accompanied by a corresponding increase in energy use.”
\end{minipage}\par
\vspace{10mm}
\textbf{Date:} September 2nd 2026\quad\textbullet\quad \textbf{Level:} B1--mid-B2 TOEFL
\vfill
Created and compiled by \textbf{Amir Kooshky}\par
\vspace{2mm}
\href{https://t.me/KooshkyTOEFL}{Telegram Channel}\quad\textbullet\quad
\href{https://www.instagram.com/kooshkytoefl}{Instagram}\quad\textbullet\quad
\href{https://t.me/KooshkyTOEFL\_pv}{Personal Telegram}
\end{center}
\newpage
\tableofcontents
\newpage

\section{Meanings and usage}
\sectionrule
\subsection{Meaning 1: Directly related or matching}
\textbf{Part of speech:} adjective

\textbf{IPA:} /ˌkɔːrəˈspɑːndɪŋ/

\textbf{Pronunciation phrase:} corresponding

\textbf{Local audio file:} audios/corresponding.mp3

\textbf{Register:} neutral to formal; very common in academic, scientific, technical, and analytical writing

\textbf{Definition:} directly related to something else, especially because the two things change together, belong together, or have matching positions

\textbf{In simpler words:} matching or directly connected

\textbf{Typical pattern:} a corresponding + increase, decrease, change, number, section, or item

\subsubsection{Natural examples}
\begin{enumerate}
  \item An increase in production was accompanied by a corresponding increase in energy use.
  \item Each number on the map refers to a corresponding location in the table.
  \item Higher levels of education are often associated with a corresponding rise in average income.
  \item The decline in rainfall led to a corresponding reduction in crop yields.
  \item As the population grew, there was a corresponding increase in demand for housing.
  \item Each category in the chart has a corresponding symbol on the map.
  \item The researchers observed a corresponding change in behavior after the environmental conditions were altered.
  \item Lower transportation costs produced a corresponding decrease in the final price of the goods.
\end{enumerate}

\subsubsection{Nuance and implication}
\textbf{Central nuance:} Corresponding emphasizes a direct relationship or match between two things rather than a general similarity.

\textbf{Usual implication:} When one thing changes, appears, or occupies a particular position, the corresponding thing is the one directly connected with it.

\textbf{Compared with unrelated:} Corresponding things have a clear relationship or match, while unrelated things have no such direct connection.

\subsubsection{Grammar notes}
\textbf{How to use it:} Use corresponding before a noun when identifying the thing that directly matches or relates to another thing.

\textbf{What can follow it:} Common nouns include increase, decrease, change, rise, fall, number, section, category, value, and item.

\textbf{Common preposition:} When the other thing must be stated explicitly, corresponding to is common, as in the value corresponding to each category.

\textbf{Position in a sentence:} Corresponding is especially common before nouns, as in a corresponding increase or the corresponding section.

\subsubsection{Useful patterns}
\textbf{Pattern:} a corresponding increase or rise

\textbf{Use:} Use this when one increase is directly connected with another.

\textbf{Example:} The expansion of the city produced a corresponding increase in traffic.

\textbf{Pattern:} a corresponding decrease or reduction

\textbf{Use:} Use this when a decrease is directly related to another change.

\textbf{Example:} Greater efficiency resulted in a corresponding reduction in fuel consumption.

\textbf{Pattern:} a corresponding change

\textbf{Use:} Use this for a change that matches or results from another change.

\textbf{Example:} The shift in temperature caused a corresponding change in the animals' behavior.

\textbf{Pattern:} the corresponding + item, number, section, or category

\textbf{Use:} Use this to identify the item that matches another item in a different list, table, diagram, or system.

\textbf{Example:} Match each definition with the corresponding term.

\textbf{Pattern:} corresponding to + noun

\textbf{Use:} Use this when explicitly identifying what something matches.

\textbf{Example:} The values corresponding to each age group are shown in the final column.

\subsubsection{Common wrong usage}
\paragraph{Correction 1}
\textbf{Wrong:} The population increased, with a correspond increase in housing demand.

\textbf{Correct:} The population increased, with a corresponding increase in housing demand.

\textbf{Why:} Use the adjective corresponding before a noun such as increase.

\paragraph{Correction 2}
\textbf{Wrong:} Each number has a corresponding with a location on the map.

\textbf{Correct:} Each number has a corresponding location on the map.

\textbf{Why:} Corresponding can directly modify a noun; with is not needed in this structure.

\paragraph{Correction 3}
\textbf{Wrong:} The corresponding increase means the two increases are exactly equal.

\textbf{Correct:} The corresponding increase means the two changes are directly related.

\textbf{Why:} Corresponding does not necessarily mean identical in size. It means directly related or matching in some relevant way.

\subsection{Meaning 2: Equivalent in another system or group}
\textbf{Part of speech:} adjective

\textbf{IPA:} /ˌkɔːrəˈspɑːndɪŋ/

\textbf{Pronunciation phrase:} corresponding

\textbf{Local audio file:} audios/corresponding.mp3

\textbf{Register:} neutral to formal; common in academic, technical, administrative, and comparative contexts

\textbf{Definition:} having the same or a very similar function, position, or meaning as something in another system, group, or situation

\textbf{In simpler words:} the matching equivalent

\textbf{Typical pattern:} the corresponding + part, section, term, position, value, or item

\subsubsection{Natural examples}
\begin{enumerate}
  \item Researchers compared each region with the corresponding region in the earlier survey.
  \item The French term has no exact corresponding expression in English.
  \item Students should enter each answer in the corresponding box on the form.
  \item Each department sends its report to the corresponding office at the national level.
  \item The two diagrams use different colors to identify corresponding parts of the structure.
  \item The researchers compared the experimental group with the corresponding control group.
  \item Each chapter in the revised edition includes a corresponding section in the online materials.
  \item The numbered samples were placed beside their corresponding labels.
\end{enumerate}

\subsubsection{Nuance and implication}
\textbf{Central nuance:} This sense emphasizes equivalence of role, position, or meaning across two separate systems, groups, documents, or sets.

\textbf{Usual implication:} The two things are not necessarily identical, but each occupies a comparable place or performs a comparable function.

\textbf{Compared with identical:} Corresponding things match in role, position, or relationship, but they do not have to be completely identical.

\subsubsection{Grammar notes}
\textbf{How to use it:} Use corresponding before the noun that has an equivalent in another group or system.

\textbf{What can follow it:} Common nouns include term, section, position, value, part, office, category, and label.

\textbf{Common preposition:} Corresponding to can identify the thing that provides the match: the English term corresponding to the French expression.

\textbf{Position in a sentence:} The definite article the is common when a unique matching item is intended: the corresponding section.

\subsubsection{Useful patterns}
\textbf{Pattern:} the corresponding section or part

\textbf{Use:} Use this when identifying the equivalent part of another document, object, or system.

\textbf{Example:} The same information appears in the corresponding section of the second report.

\textbf{Pattern:} the corresponding term or expression

\textbf{Use:} Use this when comparing vocabulary or terminology across languages or fields.

\textbf{Example:} There is no precise corresponding term in English.

\textbf{Pattern:} the corresponding value

\textbf{Use:} Use this for a value that matches a particular variable, category, or position.

\textbf{Example:} Find the temperature and record the corresponding pressure value.

\textbf{Pattern:} corresponding parts

\textbf{Use:} Use this for parts in separate objects or diagrams that occupy equivalent positions.

\textbf{Example:} The corresponding parts of the two organisms perform similar functions.

\subsubsection{Common wrong usage}
\paragraph{Correction 1}
\textbf{Wrong:} Write your answer in the correspondingly box.

\textbf{Correct:} Write your answer in the corresponding box.

\textbf{Why:} Use the adjective corresponding before the noun box.

\paragraph{Correction 2}
\textbf{Wrong:} Corresponding means completely identical.

\textbf{Correct:} Corresponding means matching in position, function, relationship, or meaning.

\textbf{Why:} Corresponding things may differ in details even though they occupy equivalent roles.

\paragraph{Correction 3}
\textbf{Wrong:} The English word is corresponding with the French term.

\textbf{Correct:} The English word corresponds to the French term.

\textbf{Why:} When using the verb correspond to describe equivalence, correspond to is usually more natural than be corresponding with.

\section{Word family}
\sectionrule
\subsection{correspond}
\textbf{Part of speech:} intransitive verb

\textbf{IPA:} /ˌkɔːrəˈspɑːnd/

\textbf{Definition:} to match, agree with, or be closely related to something

\textbf{Relationship to the headword:} Corresponding is formed from the verb correspond and describes something that matches or is directly related to another thing.

\textbf{Grammar pattern:} correspond to or correspond with + noun

\textbf{Usage note:} Correspond to is especially common for equivalence or matching. Correspond with can also mean agree with, and it has a separate meaning of exchange letters or messages.

\subsubsection{Examples}
\begin{enumerate}
  \item The results correspond closely to the predictions of the model.
  \item Each symbol corresponds to a particular type of material.
  \item The boundaries shown on the map do not correspond exactly with modern political borders.
\end{enumerate}

\subsection{correspondence}
\textbf{Part of speech:} countable or uncountable noun

\textbf{IPA:} /ˌkɔːrəˈspɑːndəns/

\textbf{Definition:} a close similarity or relationship between two things

\textbf{Relationship to the headword:} This noun can name the matching relationship expressed by corresponding.

\textbf{Grammar pattern:} a correspondence between + two things

\textbf{Usage note:} Correspondence also commonly means letters or messages exchanged between people.

\subsubsection{Examples}
\begin{enumerate}
  \item Researchers found a close correspondence between the two sets of data.
  \item There is no exact correspondence between the categories used by the two studies.
\end{enumerate}

\subsection{correspondingly}
\textbf{Part of speech:} adverb

\textbf{IPA:} /ˌkɔːrəˈspɑːndɪŋli/

\textbf{Definition:} in a way that matches or changes in relation to something else

\textbf{Relationship to the headword:} This adverb describes an action or change that occurs in a matching or related way.

\textbf{Grammar pattern:} verb + correspondingly, or correspondingly + verb

\textbf{Usage note:} It is formal and less common than the adjective corresponding.

\subsubsection{Examples}
\begin{enumerate}
  \item Production costs fell, and prices decreased correspondingly.
  \item As demand increased, the company correspondingly expanded its workforce.
\end{enumerate}

\section{Common collocations}
\sectionrule
\subsection{corresponding increase}
\textbf{Applies to:} Directly related or matching

\textbf{Meaning:} an increase directly related to another increase or change

\textbf{Example:} The increase in population produced a corresponding increase in demand for water.

\subsection{corresponding decrease}
\textbf{Applies to:} Directly related or matching

\textbf{Meaning:} a decrease directly associated with another change

\textbf{Example:} The reduction in production led to a corresponding decrease in energy consumption.

\subsection{corresponding rise}
\textbf{Applies to:} Directly related or matching

\textbf{Meaning:} a rise that occurs in relation to another change

\textbf{Example:} Higher temperatures were accompanied by a corresponding rise in evaporation.

\subsection{corresponding reduction}
\textbf{Applies to:} Directly related or matching

\textbf{Meaning:} a reduction directly related to another development

\textbf{Example:} Improved efficiency produced a corresponding reduction in operating costs.

\subsection{corresponding change}
\textbf{Applies to:} Directly related or matching

\textbf{Meaning:} a change that matches or is directly related to another change

\textbf{Example:} The change in diet produced a corresponding change in growth rates.

\subsection{corresponding value}
\textbf{Applies to:} Equivalent in another system or group

\textbf{Meaning:} the value that matches a particular item, category, or variable

\textbf{Example:} Students recorded the temperature and its corresponding pressure value.

\subsection{corresponding section}
\textbf{Applies to:} Equivalent in another system or group

\textbf{Meaning:} the section occupying an equivalent position in another document or system

\textbf{Example:} Additional details appear in the corresponding section of the report.

\subsection{corresponding part}
\textbf{Applies to:} Equivalent in another system or group

\textbf{Meaning:} a part that occupies the same or a similar position or function

\textbf{Example:} The corresponding parts of the two structures were compared.

\subsection{corresponding term}
\textbf{Applies to:} Equivalent in another system or group

\textbf{Meaning:} a word or expression with an equivalent meaning

\textbf{Example:} The language has no exact corresponding term for the concept.

\subsection{corresponding to}
\textbf{Applies to:} all senses

\textbf{Meaning:} matching or directly related to something

\textbf{Pattern:} corresponding to + noun

\textbf{Example:} The numbers corresponding to each category are shown in the table.

\section{Synonyms and antonyms}
\sectionrule
\subsection{Close synonyms}
\subsubsection{matching}
\textbf{Part of speech:} adjective

\textbf{Applies to:} Equivalent in another system or group

\textbf{Shared meaning:} Both can describe things that belong together or have equivalent roles.

\textbf{Difference:} Matching is the more general and everyday word. Corresponding is more formal and often emphasizes equivalent position, function, or relationship rather than visual similarity.

\textbf{Example:} Place each card beside its matching label.

\subsubsection{equivalent}
\textbf{Part of speech:} adjective

\textbf{Applies to:} Equivalent in another system or group

\textbf{Shared meaning:} Both can describe things with comparable roles or meanings.

\textbf{Difference:} Equivalent emphasizes having the same value, function, meaning, or effect. Corresponding emphasizes that two things occupy matching positions or are directly related.

\textbf{Example:} The two qualifications are considered academically equivalent.

\subsubsection{related}
\textbf{Part of speech:} adjective

\textbf{Applies to:} Directly related or matching

\textbf{Shared meaning:} Both describe a connection between things.

\textbf{Difference:} Related is much broader and simply means connected in some way. Corresponding suggests a more specific one-to-one or parallel relationship.

\textbf{Example:} The researchers examined several related environmental factors.

\subsubsection{parallel}
\textbf{Part of speech:} adjective

\textbf{Applies to:} Directly related or matching

\textbf{Shared meaning:} Both can describe similar or related developments.

\textbf{Difference:} Parallel often describes similar developments occurring side by side. Corresponding more directly identifies one thing as the match or counterpart of another.

\textbf{Example:} The two countries experienced parallel economic changes.

\subsection{Direct or useful antonyms}
\subsubsection{unrelated}
\textbf{Part of speech:} adjective

\textbf{Applies to:} Directly related or matching

\textbf{Opposition:} Corresponding things have a direct relationship or match, while unrelated things have no meaningful connection.

\textbf{Example:} The two changes appear to be unrelated.

\subsubsection{nonmatching}
\textbf{Part of speech:} adjective

\textbf{Applies to:} Equivalent in another system or group

\textbf{Opposition:} Corresponding items match in position, role, or relationship, while nonmatching items do not.

\textbf{Usage note:} This contrast is useful in technical contexts but is less common than unrelated.

\textbf{Example:} The software highlighted the nonmatching entries in the two datasets.

\section{Words learners may confuse}
\sectionrule
\subsection{correspondent}
\textbf{Part of speech:} adjective versus countable noun

\textbf{Why learners confuse them:} The words come from the same verb and look very similar.

\textbf{Key difference:} Corresponding is usually an adjective meaning matching or related. A correspondent is normally a person who reports news or exchanges letters or messages with someone.

\textbf{corresponding:} Enter each number in the corresponding box.

\textbf{correspondent:} The newspaper sent a foreign correspondent to report on the conflict.

\subsection{correspond}
\textbf{Part of speech:} adjective versus intransitive verb

\textbf{Why learners confuse them:} They express the same basic relationship but have different grammatical roles.

\textbf{Key difference:} Correspond is a verb meaning match, agree, or be equivalent. Corresponding is the adjective used before a noun to describe the matching thing.

\textbf{corresponding:} Each diagram has a corresponding explanation.

\textbf{correspond:} Each diagram corresponds to an explanation in the text.

\subsection{respective}
\textbf{Part of speech:} adjective

\textbf{Why learners confuse them:} Both are often used when pairing or distinguishing items in formal writing.

\textbf{Key difference:} Respective means belonging separately to each of two or more people or things already mentioned. Corresponding means matching or occupying an equivalent position.

\textbf{corresponding:} Match each graph with its corresponding description.

\textbf{respective:} The two researchers returned to their respective universities.

\section{Etymology}
\sectionrule
\textbf{Origin:} Corresponding comes from correspond, which ultimately comes from Latin elements meaning to answer together or respond in agreement.

\textbf{Development:} The idea of things answering or matching one another developed into the modern sense of having a direct relationship or equivalent position.

\textbf{Memory link:} If two things correspond, they match each other; the corresponding item is the one that matches.

\section{Practice exercises}
\sectionrule
\subsection{Word family choice}
Choose the best member of the word family for each blank.

\begin{enumerate}
  \item The values in the second table \_\_\_\_\_ closely to those reported in the first study.
  \item Researchers found a close \_\_\_\_\_ between the two patterns.
  \item As production increased, energy consumption rose \_\_\_\_\_.
  \item Each symbol has a \_\_\_\_\_ explanation in the table.
\end{enumerate}

\subsection{Correct or incorrect?}
Decide whether each sentence is correct. Correct the inaccurate ones.

\begin{enumerate}
  \item The rise in production was accompanied by a corresponding increase in energy use.
  \item Each number has a correspondingly symbol on the map.
  \item The two corresponding parts must be completely identical in every detail.
  \item The value corresponding to each category is shown in the final column.
  \item The English term corresponds to the expression used in French.
\end{enumerate}

\subsection{Collocation practice}
Complete each sentence with a natural collocation.

\begin{enumerate}
  \item The increase in demand produced a corresponding \_\_\_\_\_ in prices.
  \item Match each diagram with the \_\_\_\_\_ description.
  \item The values \_\_\_\_\_ to each age group appear in the second column.
  \item A reduction in rainfall led to a corresponding \_\_\_\_\_ in agricultural output.
\end{enumerate}

\subsection{Complete answer key}
\subsubsection{Word family choice}
\begin{enumerate}
  \item \textbf{correspond.} The values in the second table correspond closely to those reported in the first study. \emph{Use the verb correspond when saying that two sets of values match.}
  \item \textbf{correspondence.} Researchers found a close correspondence between the two patterns. \emph{Use the noun correspondence for a close matching relationship.}
  \item \textbf{correspondingly.} As production increased, energy consumption rose correspondingly. \emph{Use the adverb correspondingly to describe a matching change.}
  \item \textbf{corresponding.} Each symbol has a corresponding explanation in the table. \emph{Use the adjective corresponding before explanation.}
\end{enumerate}

\subsubsection{Correct or incorrect?}
\begin{enumerate}
  \item \textbf{Correct} \emph{A corresponding increase is an increase directly related to another change.}
  \item \textbf{Incorrect}. Each number has a corresponding symbol on the map. \emph{Use the adjective corresponding before the noun symbol.}
  \item \textbf{Incorrect}. The two corresponding parts occupy equivalent positions or have related functions. \emph{Corresponding does not require complete identity.}
  \item \textbf{Correct} \emph{Corresponding to is a natural pattern for identifying a matching value.}
  \item \textbf{Correct} \emph{Correspond to is commonly used when two things are equivalent or match.}
\end{enumerate}

\subsubsection{Collocation practice}
\begin{enumerate}
  \item \textbf{rise.} The increase in demand produced a corresponding rise in prices. \emph{A corresponding rise is a rise directly related to another change.}
  \item \textbf{corresponding.} Match each diagram with the corresponding description. \emph{The corresponding description is the one that matches a particular diagram.}
  \item \textbf{corresponding.} The values corresponding to each age group appear in the second column. \emph{Corresponding to identifies the values associated with each group.}
  \item \textbf{decrease.} A reduction in rainfall led to a corresponding decrease in agricultural output. \emph{A corresponding decrease is a decrease directly associated with another change.}
\end{enumerate}

\vfill
\begin{center}
\small\color{Muted} Created and compiled by \textbf{Amir Kooshky}\\
\href{https://t.me/KooshkyTOEFL}{Telegram Channel}\quad\textbullet\quad
\href{https://www.instagram.com/kooshkytoefl}{Instagram}\quad\textbullet\quad
\href{https://t.me/KooshkyTOEFL\_pv}{Personal Telegram}
\end{center}
\end{document}
